\section{Introduction}

Multi-agent systems built on recursive spawning architectures are susceptible to a failure mode we term \textbf{autonomous drift}. As agents spawn child agents — and those children spawn further descendants — the lineage graph grows unbounded. Without structural constraints, descendants can diverge from the originating intent, latch onto unrelated objectives, or accumulate state that silently severs the semantic thread connecting them to their ancestry. The result is an agent lineage that continues to operate, consume resources, and produce outputs, yet no longer meaningfully serves the purpose for which it was originally instantiated. This failure mode is particularly insidious because it emerges gradually: no single spawn event is obviously wrong, and the drift only becomes apparent when a cascade of descendant agents produces outcomes that are contextually misaligned with the root objective.

Existing approaches to mitigating drift fall into two broad categories, both of which prove insufficient under scrutiny. The first category is \textbf{time-to-live (TTL) constraints}, which bound the operational lifespan of any spawned agent. TTL-based systems terminate agents after a fixed wall-clock duration or a bounded number of action cycles. While TTL imposes a hard ceiling on resource consumption, it does nothing to constrain \textit{what} an agent does during its window of operation. A drifted descendant operating within its TTL window can cause significant harm — generating misaligned outputs, writing corrupted state, or initiating downstream spawns — before being forcibly terminated. Moreover, TTL-based systems offer no principled way to set the duration parameter: too short and legitimate long-horizon tasks are aborted; too long and drift has ample time to propagate. The TTL approach treats the symptom (unbounded operation) rather than the cause (unconstrained lineage divergence).

The second category is \textbf{depth-limited spawning}, which caps the number of generational levels a lineage can contain. Under depth limits, the root agent may spawn children, but those children may not spawn further descendants. Depth limits succeed at flattening the lineage graph and preventing exponential spawn cascades, but they do so at the cost of expressiveness. Many legitimate multi-agent workflows require multi-hop delegation: a planning agent spawns a research agent, which spawns a synthesis agent, which spawns a citation agent. Depth limits render such workflows impossible or force artificial restructuring around a flat hierarchy. More critically, depth limits do not prevent drift within a single generation — a depth-limited child can still diverge from its parent's intent and propagate that divergence through any side-channels (file system state, shared memory, external API calls) that remain accessible within the depth cap. Depth is a proxy for complexity, not for alignment.

\subsection{Core Insight: Immutable Parent Pointer Chain}

The structural failure underlying autonomous drift is the absence of an immutable, verifiable lineage anchor. In systems that permit dynamic re-parenting, speculative spawning with deferred commitment, or lineage graph mutations after the fact, descendant agents lack a stable reference point against which to validate continued alignment. Drift is possible precisely because there is no immutable record of the chain of intentions that produced a given agent.

We observe that drift becomes \textit{structurally impossible} when every agent is instantiated with an immutable parent pointer that forms a cryptographically tamper-evident chain back to a genesis event. In the BX3 framework, this chain is implemented through the \textbf{Substrate}, which records each spawn event as an immutable entry containing: the parent agent's canonical identifier, a content hash of the spawning context, a timestamp, and the spawner's cryptographic signature. The entry is written to the Substrate's append-only event log and cannot be modified retroactively. Consequently, any descendant agent can be reconstructed — by any observer, at any time — as a path through the parent pointer chain from genesis to the present agent.

The immutability of this chain provides two properties that jointly eliminate drift. First, \textbf{causal traceability}: because each agent's genesis can be traced unambiguously to a genesis event, there is no ambiguity about the provenance of any agent in the system. Second, \textbf{semantic continuity}: the chain of spawning contexts constitutes a cryptographically verifiable record of the entire objective propagation path. An agent that deviates from its declared intent can be detected by traversing the chain and comparing each spawning context against the agent's actual behavior — a discrepancy between context and behavior indicates drift, and the chain provides the evidence trail for remediation.

Formally, we define \textbf{drift} as follows. Let $A_0$ denote a genesis agent instantiated at time $t_0$ with an explicit objective $o_0$. For any agent $A_i$ spawned at time $t_i$, let $p(A_i)$ denote its parent agent and let $c_i$ denote the spawning context recorded in the Substrate at spawn time. An agent $A_n$ is said to have \textit{drifted} at time $t > t_n$ if there exists a path $A_0, A_1, \ldots, A_n$ such that for all $k \in \{1, \ldots, n\}$, $p(A_k) = A_{k-1}$ and the objective $o_k$ encoded in $c_k$ is not semantically subsumed by the objective $o_{k-1}$ encoded in $c_{k-1}$, and where $A_n$'s actual behavior at time $t$ is not directed toward any objective in the transitive closure of $\{o_0, o_1, \ldots, o_n\}$. Drift is thus a property of semantic divergence between the recorded lineage and the agent's operational output — it is detectable by traversing the parent pointer chain, reconstructing the transitive objective closure, and comparing it against observable agent behavior.

Under the immutable parent pointer architecture, drift cannot occur without a cryptographic break in the chain (which the Substrate prevents) or a misrepresentation of objective propagation at spawn time (which the content-hash integrity check prevents). The chain thus transforms drift from a systemic risk requiring continuous monitoring into a structurally eliminated failure mode.

\subsection{Paper Roadmap}

This paper is organized as follows. Section 2 presents the formal model of recursive spawning under the immutable parent pointer architecture, establishing the theoretical underpinnings of the BX3 agentic substrate. Section 3 describes the Substrate implementation, including the append-only event log, content-hash integrity mechanism, and the spawn authorization protocol. Section 4 analyzes the drift-elimination properties of the architecture through a series of structured agentic workflows, demonstrating that drift is not merely reduced but structurally impossible under the stated constraints. Section 5 presents empirical evaluation of multi-agent spawn chains operating under the BX3 substrate, measuring lineage integrity, objective alignment retention, and cascade backstop trigger accuracy. Section 6 discusses related work in multi-agent safety, lineage tracking, and autonomous system alignment. Section 7 concludes with a discussion of open questions and planned extensions to the BX3 substrate, including extension of the immutable pointer chain to support horizontal agent collaboration and cross-chain spawning attestation.